\section{\SystemName{} 在生产中的实践}
\label{sec:large-scale-analysis}

过去六个月里，已有数千个 agentic RL 作业使用 \SystemName{} 进行后训练。为进一步展示系统的可扩展性与鲁棒性，我们报告其在一个超过 3000 张 GPU 的生产集群上的部署经验。

\PHM{工作负载画像}
我们使用 \SystemName{} 在内部数学与软件工程数据集上训练了一个百亿级以上参数规模的 MoE LLM。该任务的最大 prompt 长度和回复长度分别为 12k 与 46k，单任务平均交互轮次从 1 到 48 不等（见 \autoref{fig:action}），这再次表明 agentic RL 中 prefill 主导与 decode 主导任务会同时存在。大量轮次意味着环境必须极为稳定，同时 prefill 计算也要足够快。

该作业采用异步训练，训练 GPU 与生成 GPU 的比例为 1:5。为了兼顾梯度稳定性与 rollout 效率，我们将异步边界设为 1，对应的最大单步时长约为 1.5 小时。系统的主要瓶颈来自阻塞式 \texttt{get\_batch} 调用：训练阶段完成 log probability 与梯度计算后，仍需等待 \texttt{SampleBuffer} 中积累足够多的轨迹。这一等待最高可占迭代时间的 62\%，若能完全消除，理想情况下端到端训练时间可再下降 22\%。

\PHM{基于画像的优化}
根据上述工作负载分析，我们调整了训练与 rollout 的资源比例，并针对 MoE 架构优化了 prefix caching。\autoref{fig:opt_step_time} 表明，这些画像驱动优化在前 25 个 step 内带来了 1.66 倍的端到端加速。除此之外，系统还对环境稳定性与故障恢复进行了专门优化。

\PHM{环境稳定性优化}
在 Kubernetes 上同时运行数千个基于 Docker 的 agentic 环境，对稳定性和性能都提出了较高要求。为减少 \texttt{env.reset} 中的网络抖动与镜像拉取开销，我们部署了多级缓存架构：第一层为内部镜像仓库镜像源，避免外网访问；第二层为位于计算节点与仓库之间的分布式负载均衡缓存，用于应对大规模并发拉取请求。在 \texttt{env.reset} 过程中，客户端优先从缓存获取 Docker 镜像。该方案将环境初始化成功率提升到 99.99\% 以上。

\noindent\textbf{系统韧性。}
我们同时从网络与系统两个层面增强鲁棒性。为减少连接环境时的超时，系统采用基于会话的持久协议而非一次一请求协议，并用指数退避处理临时网络错误。解耦式架构还能隔离故障，使单个环境、奖励 worker 或推理 worker 的问题不至于扩散到其他部分。环境 worker 由 Kubernetes 管理，奖励 worker 由 serverless 平台管理；推理 worker 若失败，首先在原 GPU 上重启，再失败则被剔除，其轨迹从存储引擎恢复并转移到健康 worker；训练 worker 若失败，则从最近检查点重启。实际生产运行一周仅观察到一次故障，证明了这一设计的稳健性。
